\documentclass[11pt]{article}
\usepackage[margin=1in]{geometry}
\usepackage{amsmath,amsthm,amssymb,mathtools}
\usepackage[colorlinks=true,linkcolor=blue,citecolor=blue,urlcolor=blue]{hyperref}

\newtheorem{theorem}{Theorem}
\newtheorem{corollary}[theorem]{Corollary}
\newtheorem{lemma}[theorem]{Lemma}
\theoremstyle{remark}
\newtheorem{remark}[theorem]{Remark}

\newcommand{\Cay}{\operatorname{Cay}}
\newcommand{\Cdir}{\vec C}
\newcommand{\torus}[2]{D_{#1}(#2)}

\title{A Cartesian-power lift for Hamilton decompositions of directed tori}
\author{}
\date{}

\begin{document}
\maketitle

\begin{abstract}
We isolate a general reduction behind composite-dimension constructions for directed tori.
If a finite digraph $G$ admits a Hamilton decomposition and the directed torus
$\vec C_{|V(G)|}^{\square b}$ also admits one, then the Cartesian power $G^{\square b}$ admits a
Hamilton decomposition.
Applied to
\[
\torus{d}{m}=\Cay\bigl((\mathbb Z/m\mathbb Z)^d,\{e_0,\dots,e_{d-1}\}\bigr),
\]
this yields the pointwise implication
\[
\torus{a}{m}\text{ and }\torus{b}{m^a}\text{ decomposable }\Longrightarrow \torus{ab}{m}\text{ decomposable},
\]
and therefore a multiplicative-closure reduction for any family of base dimensions already known uniformly in the modulus.
The contribution of this note is thus theorem-level rather than base-case level: it packages a clean product argument that propagates solved dimensions to composite ones.
\end{abstract}

\section{Literature context and scope}
Cartesian products of directed cycles have a substantial Hamiltonicity literature.
Trotter and Erd\H{o}s characterized when $\vec C_{n_1}\square \vec C_{n_2}$ is Hamiltonian \cite{TrotterErdos1978}.
Curran and Witte characterized Hamilton paths in $\vec C_{n_1}\square \vec C_{n_2}$ and showed that every Cartesian product of three or more nontrivial directed cycles has a Hamilton cycle \cite{CurranWitte1985}; see also the surveys \cite{CurranGallian1996,LanelPallageRatnayakeThevashaWelihinda2019}.
More recently, Bogdanowicz proved a common-factor criterion for decompositions of Cartesian products of directed cycles into equal-length directed cycles \cite{Bogdanowicz2017}, and identified Hamilton cycles and, in some arithmetic cases, two arc-disjoint Hamilton cycles in products of directed cycles \cite{Bogdanowicz2020}.
Darijani, Miraftab, and Morris proved that products of two directed cycles, and more generally of any four or more directed cycles, admit two arc-disjoint Hamiltonian paths \cite{DarijaniMiraftabMorris2025}.

In the undirected setting, Hamilton decompositions of Cartesian products have long been studied; classical product-decomposition results include Foregger \cite{Foregger1978}, Aubert--Schneider \cite{AubertSchneider1982}, and Stong \cite{Stong1991}.
Within the equal-side directed-torus program relevant here, companion manuscript-level work includes a complete $d=3$ paper for all $m\ge 3$ and a finished $d=4$ proof program recorded in Park's repository bundle \cite{ParkD3,ParkRepo2026}.
Against this background, the present note is narrower but different.
It does not attempt to settle new base dimensions, and it does not address arbitrary mixed products $\vec C_{n_1}\square\cdots\square\vec C_{n_r}$.
Instead, it isolates a Cartesian-power lift tailored to equal-side directed tori: once dimensions $a$ and $b$ are available uniformly in the modulus, the theorem here propagates them to the composite dimension $ab$.

\section{General Cartesian-power lift}

Let $G=(V,E)$ be a finite digraph on $n$ vertices.
Recall that a Hamilton decomposition of $G$ is a partition of $E$ into directed Hamilton cycles.
Write $G^{\square b}$ for the $b$-fold Cartesian power of $G$.

\begin{theorem}[Cartesian-power lift]\label{thm:product}
Let $G$ be a finite digraph on $n$ vertices.
Assume:
\begin{enumerate}
    \item $G$ admits a Hamilton decomposition into $a$ directed Hamilton cycles;
    \item $\vec C_n^{\square b}$ admits a Hamilton decomposition into $b$ directed Hamilton cycles.
\end{enumerate}
Then $G^{\square b}$ admits a Hamilton decomposition into $ab$ directed Hamilton cycles.
\end{theorem}

\begin{proof}
Choose a Hamilton decomposition
\[
E(G)=H_1\sqcup H_2\sqcup\cdots\sqcup H_a,
\]
where each $H_i$ is a directed Hamilton cycle on $V$.

Fix $i\in\{1,\dots,a\}$.
In $G^{\square b}$, consider the spanning subdigraph
\[
K_i:=H_i^{\square b}.
\]
Since each $H_i$ is a directed $n$-cycle, we have a canonical isomorphism
\[
K_i\cong \vec C_n^{\square b}.
\]
By assumption, $\vec C_n^{\square b}$ has a Hamilton decomposition into $b$ directed Hamilton cycles, so $K_i$ also has one.
Write
\[
E(K_i)=A_i^{(1)}\sqcup A_i^{(2)}\sqcup\cdots\sqcup A_i^{(b)},
\]
where each $A_i^{(r)}$ is a directed Hamilton cycle on $V^b$.

Now every arc of $G^{\square b}$ belongs to exactly one of the subdigraphs $K_i$:
an arc changes exactly one coordinate, and in that coordinate it uses exactly one arc of $G$, which lies in exactly one Hamilton factor $H_i$.
Hence
\[
E(G^{\square b})=E(K_1)\sqcup E(K_2)\sqcup\cdots\sqcup E(K_a).
\]
Substituting the decompositions of the $K_i$, we obtain
\[
E(G^{\square b})=\bigsqcup_{i=1}^a\bigsqcup_{r=1}^b A_i^{(r)},
\]
a partition into $ab$ directed Hamilton cycles.
So $G^{\square b}$ admits a Hamilton decomposition.
\end{proof}

\begin{corollary}[Self-square closure]\label{cor:selfsquare}
If a finite digraph $G$ admits a Hamilton decomposition, then $G\square G$ admits a Hamilton decomposition.
\end{corollary}

\begin{proof}
Take $b=2$ in Theorem~\ref{thm:product}.
Since $\vec C_n\square\vec C_n$ is a 2-regular directed torus, it suffices to know that it decomposes into two directed Hamilton cycles; an explicit proof appears in Lemma~\ref{lem:square} below.
\end{proof}

\section{Application to directed tori}

For integers $d\ge 1$ and $m\ge 2$, define the directed $d$-torus
\[
\torus{d}{m}=\Cay\bigl((\mathbb Z/m\mathbb Z)^d,\{e_0,\dots,e_{d-1}\}\bigr).
\]
Equivalently, this is the Cartesian power $\vec C_m^{\square d}$.

\begin{theorem}[Pointwise composite expansion]\label{thm:abm}
Let $a,b\ge 1$ and $m\ge 2$.
If both $\torus{a}{m}$ and $\torus{b}{m^a}$ admit Hamilton decompositions, then $\torus{ab}{m}$ admits a Hamilton decomposition.
\end{theorem}

\begin{proof}
Observe that
\[
\torus{ab}{m}=\bigl(\vec C_m^{\square a}\bigr)^{\square b}=\bigl(\torus{a}{m}\bigr)^{\square b}.
\]
Apply Theorem~\ref{thm:product} with $G=\torus{a}{m}$.
The digraph $G$ has $m^a$ vertices, so the auxiliary input required by Theorem~\ref{thm:product} is exactly
\[
\vec C_{m^a}^{\square b}=\torus{b}{m^a}.
\]
By hypothesis both inputs admit Hamilton decompositions, hence so does $G^{\square b}=\torus{ab}{m}$.
\end{proof}

\begin{remark}
Theorem~\ref{thm:abm} is deliberately pointwise in the modulus: to pass from dimensions $a$ and $b$ to the composite dimension $ab$, the second input is $\torus{b}{m^a}$ rather than $\torus{b}{m}$.
The uniform multiplicative-closure corollary below is obtained only after quantifying this pointwise implication over all moduli.
\end{remark}

\begin{corollary}[Uniform multiplicative closure]\label{cor:multiplicative}
Suppose $A\subseteq \{1,2,3,\dots\}$ has the property that for every $d\in A$ and every integer $m\ge 3$, the directed torus $\torus{d}{m}$ admits a Hamilton decomposition.
Then the same is true for every dimension in the multiplicative semigroup generated by $A$.
\end{corollary}

\begin{proof}
It suffices to prove closure under multiplication.
Take $a,b\in A$ and fix $m\ge 3$.
By hypothesis, $\torus{a}{m}$ admits a Hamilton decomposition, and so does
\[
\torus{b}{m^a}
\]
because $m^a\ge 3$.
Therefore Theorem~\ref{thm:abm} implies that $\torus{ab}{m}$ admits a Hamilton decomposition.
\end{proof}

\section{The square case d=2}

\begin{lemma}[Two-cycle decomposition of the square torus]\label{lem:square}
For every integer $n\ge 2$, the directed torus
\[
\vec C_n\square \vec C_n
\]
decomposes into two directed Hamilton cycles.
\end{lemma}

\begin{proof}
Label vertices by pairs $(i,j)\in \mathbb Z_n^2$.
From each vertex there are two outgoing arcs:
\[(i,j)\to(i+1,j),\qquad (i,j)\to(i,j+1).
\]
Define two 1-factors $F_0,F_1$ by
\[
F_0(i,j)=
\begin{cases}
(i+1,j), & i+j\not\equiv n-1 \pmod n,\\
(i,j+1), & i+j\equiv n-1 \pmod n,
\end{cases}
\qquad
F_1(i,j)=
\begin{cases}
(i,j+1), & i+j\not\equiv n-1 \pmod n,\\
(i+1,j), & i+j\equiv n-1 \pmod n.
\end{cases}
\]
At each vertex, the two outgoing arcs are exactly partitioned between $F_0$ and $F_1$.
So it remains to show that each $F_t$ is a single directed Hamilton cycle.

For $F_0$, consider the bijection
\[
\Psi(i,j)=(i+j,\,j)
\]
with inverse $\Psi^{-1}(s,t)=(s-t,\,t)$.
A direct substitution shows
\[
\Psi F_0 \Psi^{-1}(s,t)=
\begin{cases}
(s+1,t), & s\not\equiv n-1 \pmod n,\\
(0,t+1), & s\equiv n-1 \pmod n.
\end{cases}
\]
Thus $F_0$ is conjugate to the standard odometer on $\mathbb Z_n^2$.
Starting from $(0,0)$, that odometer runs through
\[
(0,0), (1,0), \dots, (n-1,0), (0,1), \dots, (n-1,n-1)
\]
as a single orbit of length $n^2$.
Hence $F_0$ is a directed Hamilton cycle.

Let $\sigma(i,j)=(j,i)$.
Then $F_1=\sigma\circ F_0\circ\sigma^{-1}$, so $F_1$ is conjugate to $F_0$ and is also a directed Hamilton cycle.
Therefore the two 1-factors partition the arc set into two directed Hamilton cycles.
\end{proof}

\begin{remark}
The Hamiltonicity of $\vec C_{n_1}\square\vec C_{n_2}$ is classical \cite{TrotterErdos1978}.
We record the explicit equal-side factorization above because the Cartesian-power lift requires a full Hamilton decomposition of the auxiliary square torus, not merely a single Hamilton cycle.
\end{remark}

\section{Consequences for composite dimensions}

\begin{corollary}[Low-dimensional base cases and propagated dimensions]\label{cor:examples}
The preceding results have the following consequences.
\begin{enumerate}
    \item Since $\torus{2}{m}$ is solved for every $m\ge 2$ by Lemma~\ref{lem:square}, any dimension already known uniformly in the modulus automatically yields its doubled dimension.
    \item Companion work supplies uniform base cases in dimensions $3$ and $4$ for every $m\ge 3$ \cite{ParkD3,ParkRepo2026}.
    Therefore the multiplicative semigroup generated by $\{2,3,4\}$, equivalently by $\{2,3\}$, is solved for every $m\ge 3$.
    In particular, every dimension of the form $2^\alpha 3^\beta$ with $\alpha+\beta\ge 1$ admits a Hamilton decomposition for all $m\ge 3$.
    \item More generally, once uniform base cases are available for any set of prime dimensions, every composite dimension in the semigroup they generate follows formally from Corollary~\ref{cor:multiplicative}.
\end{enumerate}
\end{corollary}

\begin{proof}
Item~(1) is immediate from Corollary~\ref{cor:selfsquare} or, equivalently, from Theorem~\ref{thm:abm} with $a=d$, $b=2$.
Item~(2) is Corollary~\ref{cor:multiplicative} applied to $A=\{2,3,4\}$.
Item~(3) is exactly the content of Corollary~\ref{cor:multiplicative}.
\end{proof}

\begin{thebibliography}{99}

\bibitem{AubertSchneider1982}
Jacques Aubert and Bernadette Schneider,
\emph{Decomposition de la somme cartesienne d'un cycle et de l'union de deux cycles hamiltoniens en cycles hamiltoniens},
Discrete Mathematics \textbf{38} (1982), 7--16.

\bibitem{Bogdanowicz2017}
Zbigniew R. Bogdanowicz,
\emph{On decomposition of the Cartesian product of directed cycles into cycles of equal lengths},
Discrete Applied Mathematics \textbf{229} (2017), 148--150.

\bibitem{Bogdanowicz2020}
Zbigniew R. Bogdanowicz,
\emph{Identifying Hamilton cycles in the Cartesian product of directed cycles},
AKCE International Journal of Graphs and Combinatorics \textbf{17} (2020), 534--538.

\bibitem{CurranGallian1996}
Stephen J. Curran and Joseph A. Gallian,
\emph{Hamiltonian cycles and paths in Cayley graphs and digraphs---a survey},
Discrete Mathematics \textbf{156} (1996), 1--18.

\bibitem{CurranWitte1985}
Stephen J. Curran and David Witte,
\emph{Hamilton Paths in Cartesian Products of Directed Cycles},
Annals of Discrete Mathematics \textbf{27} (1985), 35--74.

\bibitem{DarijaniMiraftabMorris2025}
Iren Darijani, Babak Miraftab, and Dave Witte Morris,
\emph{Arc-disjoint hamiltonian paths in Cartesian products of directed cycles},
Ars Mathematica Contemporanea \textbf{25} (2025), Article P2.10.

\bibitem{Foregger1978}
Marsha F. Foregger,
\emph{Hamiltonian decompositions of products of cycles},
Discrete Mathematics \textbf{24} (1978), 251--260.

\bibitem{LanelPallageRatnayakeThevashaWelihinda2019}
G. H. J. Lanel, H. K. Pallage, J. K. Ratnayake, S. Thevasha, and B. A. K. Welihinda,
\emph{A survey on Hamiltonicity in Cayley graphs and digraphs on different groups},
Discrete Mathematics, Algorithms and Applications \textbf{11} (2019), 1930002.

\bibitem{ParkD3}
Sanghyun Park,
\emph{Hamilton decompositions of the directed 3-torus $D(m)$: explicit constructions for all $m \ge 3$},
manuscript, 2026.
Available through the repository bundle at \url{https://github.com/aria1th/Torus-Hamilton-Decomposition}.

\bibitem{ParkRepo2026}
Sanghyun Park,
\emph{Hamilton decompositions of low-dimensional directed tori},
GitHub repository, 2026.
The repository README records a solved $d=3$ manuscript, a finished $d=4$ proof program, and the composite-dimension reduction note discussed here.
\url{https://github.com/aria1th/Torus-Hamilton-Decomposition}.

\bibitem{Stong1991}
Richard Stong,
\emph{Hamilton decompositions of cartesian products of graphs},
Discrete Mathematics \textbf{90} (1991), 169--190.

\bibitem{TrotterErdos1978}
William T. Trotter, Jr. and Paul Erd\H{o}s,
\emph{When the Cartesian product of directed cycles is Hamiltonian},
Journal of Graph Theory \textbf{2} (1978), 137--142.

\end{thebibliography}

\end{document}
